% Include this section within Preliminary Experiments, after the general
% methodology and before the individual experiment results.
% Historical code: 9c676a2c, a46cc6a0, and c11f9a1a respectively.
% Run evidence: aws-results/<run-id>/work and the per-trial MLflow artifacts.
\subsection{Detailed Methodology for the Three AWS Experiments}
\label{sec:preliminary-aws-methodology}

The pressure-feature, magnetometer feature, and calibration fraction experiments
were executed from three different repository commits, as shown in
Table~\ref{tab:preliminary-methodology-runs}. Each configuration combined
NOR-TMD training windows with eligible Sydney collector windows. The Sydney
collector sessions were divided into calibration and holdout sets before model
selection. The calibration sessions contributed to training and cross-validation;
the holdout sessions were used only for the final evaluation of the selected
model. The unit of assignment was a recording group, identified by participant,
device, and session, rather than an individual window. Thus, overlapping windows
from a session could not cross the calibration/holdout boundary or a
cross-validation fold boundary. Different sessions from the same participant
could occur on both sides of a boundary.

\begin{table}[htbp]
\centering
\small
\caption{Historical code and run records used for the three methods.}
\label{tab:preliminary-methodology-runs}
\begin{tabular}{lll}
\hline
\textbf{Experiment} & \textbf{AWS run ID} & \textbf{Git commit} \\
\hline
Pressure feature & \texttt{20260809T055942Z-9c676a2c-full} & \texttt{9c676a2c} \\
Magnetometer feature & \texttt{20260812T125950Z-a46cc6a0-full} & \texttt{a46cc6a0} \\
Calibration fraction & \texttt{20260814T114054Z-c11f9a1a-full} & \texttt{c11f9a1a} \\
\hline
\end{tabular}
\end{table}

\subsubsection{Shared data preparation and model selection}

For each configured transport mode, the pipeline ingested NOR-TMD records and
Sydney collector sessions, rejected sessions shorter than 30 seconds or longer
than 28,800 seconds, and constructed windows separately within each session.
Windows were 10, 20, 30, or 60 seconds long in the first two experiments, with
steps of 5, 10, 15, or 30 seconds respectively. The calibration fraction
experiment used only 60-second windows with a 30-second step. Accelerometer,
gyroscope, and magnetometer vectors were reduced to Euclidean magnitudes before
aggregation. Pressure was treated as a scalar. All configurations used the
mean, population standard deviation, range, and minimum of accelerometer and
gyroscope magnitudes. The experiment-specific magnetometer and pressure
aggregations are stated below.

For Sydney windows, each configured sensor had to supply at least the number of
valid samples implied by its minimum sampling rate: 30~Hz for accelerometer,
4~Hz for gyroscope, and 2~Hz for magnetometer or pressure. The largest gap
between valid samples, including the window boundaries, could not exceed
500~ms. A window failing either check for any configured sensor was discarded.
NOR-TMD windows required at least one valid sample per configured sensor and
did not use these Sydney sampling and continuity checks. Consequently, adding
pressure could change which Sydney sessions yielded usable windows.

Within each configuration, five shuffled \texttt{GroupKFold} folds were formed
from the Sydney calibration groups with random seed 42. For every fold, a model
was fitted using all NOR-TMD windows plus the Sydney calibration windows in
the other four folds; only the excluded Sydney calibration windows supplied
validation predictions. Optuna used a seeded Tree-structured Parzen Estimator
to search random forest, XGBoost, and multilayer perceptron configurations,
with a budget of 45 trials per configuration. A Hyperband pruner could stop a
trial after an intermediate fold result. The objective for a completed trial
was macro F1 computed from the pooled out-of-fold Sydney predictions, rather
than the arithmetic mean of the five fold scores. Training used inverse-class-
frequency sample weights where the classifier accepted them. The winning
configuration was fitted again on all NOR-TMD and Sydney calibration windows;
in this final fit, Sydney calibration sample weights were additionally doubled.
Its macro F1 on the untouched Sydney holdout sessions was reported once. The
holdout score is therefore not a cross-validation score.

\subsubsection{Pressure-feature experiment}

At commit \texttt{9c676a2c}, 12 configurations crossed four window lengths
with three pressure choices. Every configuration used magnetometer standard
deviation and range. The no-pressure choice omitted pressure features. The two
pressure-enabled choices added pressure standard deviation and range, together
with either the last minus first valid pressure reading in each window
(window-start change) or the window mean minus the first valid pressure reading
in its session (session-baseline change).

Sydney groups were sorted by group identifier and split once using
\texttt{StratifiedShuffleSplit} with \texttt{train\_size=0.5} and
\texttt{random\_state=42}, stratifying by transport label. The four no-pressure
configurations had 140 effective Sydney groups, split into 70 calibration and
70 holdout groups. The eight pressure-enabled configurations had 104 effective
groups, split into 52 and 52. Although the trial label mapping included bus,
car, train, and tram, the pressure-enabled Sydney cohort contained no tram
sessions. Therefore, a pressure comparison also changes the effective cohort
and the classes represented in the Sydney holdout. The 0.5 setting refers to
groups; window counts need not be equal between partitions.

\subsubsection{Magnetometer feature experiment}

At commit \texttt{a46cc6a0}, 12 configurations crossed the same four
window/step pairs with three magnetometer feature sets: standard deviation and
range; those two plus minimum; and those three plus mean. Pressure was omitted
throughout. Bus, car, and train were the configured classes. The same one-time,
label-stratified group split used \texttt{train\_size=0.5} and seed 42. All
configurations used 152 effective Sydney groups, divided into 76 calibration
and 76 holdout groups. The actual window counts differed between partitions.

\subsubsection{Calibration fraction experiment}

At commit \texttt{c11f9a1a}, all three trials used 60-second windows stepped
by 30 seconds, no pressure, and magnetometer standard deviation, range, and
minimum. The 155 effective Sydney groups comprised 38 bus, 71 car, and 46
train groups. This commit replaced the single stratified shuffle split with a
seeded per-class procedure. Groups were sorted by identifier within each
class; a NumPy random generator seeded with 42 permuted each class in bus,
car, train order. For each class, the number assigned to calibration was the
floor of its group count times its requested fraction, bounded to leave at
least one group in each partition. The remainder formed the holdout. The
resulting allocations are shown in Table~\ref{tab:preliminary-methodology-calibration}.

\begin{table}[htbp]
\centering
\small
\caption{Calibration and holdout Sydney group counts by trial. Fractions and
counts are ordered bus, car, train.}
\label{tab:preliminary-methodology-calibration}
\begin{tabular}{lccc}
\hline
\textbf{Trial} & \textbf{Calibration fractions} & \textbf{Calibration groups} & \textbf{Holdout groups} \\
\hline
A & $0.8,\;0.4,\;0.4$ & $30,\;28,\;18$ & $8,\;43,\;28$ \\
B & $0.4,\;0.8,\;0.4$ & $15,\;56,\;18$ & $23,\;15,\;28$ \\
C & $0.4,\;0.4,\;0.8$ & $15,\;28,\;36$ & $23,\;43,\;10$ \\
\hline
\end{tabular}
\end{table}

\paragraph{Reproduction record.}
For each ID in Table~\ref{tab:preliminary-methodology-runs}, the matching
\texttt{aws-results} directory contains \texttt{run/run-summary.json} with
the executed Git commit, and \texttt{work} contains trial configurations and
split manifests. The \texttt{mlflow/mlartifacts} directory retains per-trial
copies of the split manifests and metrics. These records contain the actual
configuration, feature-frame fingerprint, row indices, and five fold assignments
for each trial. In the first two runs, a fixed seed reproduces the split only
when the same eligible sessions and ordered feature rows are supplied. Neither
run captured an immutable raw collector payload snapshot at execution time;
their effective collector cohorts were reconstructed from the available run
records. The calibration fraction run recorded a collector membership snapshot,
but reproduction still requires the same source data and preprocessing. Its
three trials reused one feature-work directory because their feature settings
were identical, so the final files in that directory represent only the last
trial. The separate per-trial MLflow artifacts preserve each trial's split and
metrics and are the authoritative records for the counts above.
